\documentclass[10pt]{article}
\usepackage[T1]{fontenc}
\usepackage[utf8]{inputenc}
\usepackage[french]{babel}
\usepackage{amsmath,amssymb,mathrsfs,mathpazo}
\usepackage{geometry}
\geometry{a4paper,margin=2cm}
\usepackage{graphicx,array,tabularx,tabularray,colortbl,multirow,multicol,booktabs,enumitem,pifont,ulem,fancybox,diagbox,cancel}
\usepackage{pstricks,pstricks-add,pst-plot,pst-tree,pst-node,pst-text,pst-func,pst-math,pst-bspline,pst-3dplot,pst-eucl,pst-solides3d}
\let\pscalculate\undefined
\usepackage{pst-bezier}
\usepackage{tikz,pgfplots,tkz-tab}
\usetikzlibrary{arrows,arrows.meta,patterns,calc,positioning,trees,decorations.markings,shapes.geometric,shapes.misc,shapes.arrows,intersections,decorations.pathmorphing,decorations.pathreplacing,backgrounds,shadows,babel}
\pgfplotsset{compat=1.18}
\usepgfplotslibrary{fillbetween}
\usepackage[np]{numprint}
\usepackage{listings,makecell,pas-tableur,twoopt}
\let\barre\undefined
\usepackage[pstricks]{tablvar}
\definecolor{darkgreen}{RGB}{0,100,0}
\providecommand{\red}{\color{red}}
\providecommand{\blue}{\color{blue}}
\providecommand{\green}{\color{green}}
\providecommand{\black}{\color{black}}
\providecommand{\white}{\color{white}}
\providecommand{\gray}{\color{gray}}
\providecommand{\cyan}{\color{cyan}}
\providecommand{\magenta}{\color{magenta}}
\providecommand{\yellow}{\color{yellow}}
\providecommand{\violet}{\color{violet}}
\providecommand{\orange}{\color{orange}}
\providecommand{\euro}{\texteuro}
\providecommand{\EUR}{\texteuro}
\providecommand{\ieme}{\textsuperscript{e}}
\providecommand{\degres}{\ensuremath{^\circ}}
\providecommand{\degree}{\ensuremath{^\circ}}
\providecommand{\vv}[1]{\overrightarrow{#1}}
\providecommand{\Coord}[1]{#1}
\providecommand{\decosix}{}
\providecommand{\starredbullet}{\ensuremath{\bullet}}
\pagestyle{empty}
\setlength{\parindent}{0pt}
\usepackage{adjustbox}\DeclareUnicodeCharacter{2236}{:}\DeclareUnicodeCharacter{2190}{\ensuremath{\leftarrow}}
\providecommand{\R}{\mathbb{R}}
\providecommand{\N}{\mathbb{N}}
\providecommand{\Z}{\mathbb{Z}}
\providecommand{\Q}{\mathbb{Q}}
\providecommand{\D}{\mathbb{D}}
\providecommand{\C}{\mathbb{C}}
\providecommand{\e}{{\mathrm{e}}}
\providecommand{\vect}[1]{\overrightarrow{#1}}
\providecommand{\vectt}[1]{\overrightarrow{\mathrm{#1}}}
\providecommand{\vvt}[1]{\overrightarrow{\mathrm{#1}}}
\providecommand{\np}[2][]{{#2~#1}}
\providecommand{\vv}[1]{{\overrightarrow{#1}}}
\renewcommand{\d}{{\mathrm{d}}}
\providecommand{\degres}{{^\circ}}
\providecommand{\degre}{{^\circ}}
\providecommand{\blue}{\color{blue}}
\providecommand{\red}{\color{red}}
\providecommand{\enskip}{\quad}
\providecommand{\llbracket}{\lbrack\!\lbrack}
\providecommand{\rrbracket}{\rbrack\!\rbrack}
\providecommand{\euro}{{\text{€}}}
\providecommand{\no}{{\text{n°}}}
\providecommand{\interval}[2]{{[#1;#2]}}
\providecommand{\intervFO}[2]{{[#1;#2[}}
\providecommand{\dotfill}{\ldots}

\providecommand{\ds}{\displaystyle}
\providecommand{\textding}[1]{#1}
\providecommand{\Oij}{\(\mathrm{O};\vec{\imath},\vec{\jmath}\)}
\providecommand{\Oijk}{\(\mathrm{O};\vec{\imath},\vec{\jmath},\vec{k}\)}
\providecommand{\Ouv}{\(\mathrm{O};\vec{u},\vec{v}\)}
\providecommand{\intervFF}[2]{\ensuremath{[#1;#2]}}
\renewcommand{\euro}{\ifmmode\text{\texteuro}\else\texteuro\fi}
\DeclareUnicodeCharacter{00B0}{\ifmmode^{\circ}\else\textdegree\fi}
\DeclareUnicodeCharacter{20AC}{\ifmmode\text{\texteuro}\else\texteuro\fi}
% Énoncés : Éducation nationale. Transcription et corrigés : APMEP, Denis Vergès et contributeurs. Découpage et classement : CoopMaths, Sébastien Lozano, Éric Elter et contributeurs. Adaptation de la présentation : Maths Explorer.
% CC BY-SA : https://coopmaths.fr/www/about/
% https://coopmaths.fr/alea/static/eam/2025/tex/eam_specialite_2025_06_sujet0v2_3.tex
\begin{document}
\noindent\textit{APMEP / CoopMaths. CC BY-SA. Présentation adaptée par Maths Explorer.}\par\medskip
\begin{center}
Le plan est muni d'un repère orthogonal.
\end{center}

\subsubsection*{Partie A}
On considère la fonction $P$ définie sur l'intervalle $[-5~;~3]$ par :
\[P(x) = 2x^2 + x - 10.\]

\begin{enumerate}
\item 
	\begin{enumerate}
	\item Déterminer les racines de $P$.
	\item En déduire l'axe de symétrie de la parabole d'équation $y = P(x)$.
\end{enumerate}

\item Établir le tableau de signe de la fonction $P$ sur l'intervalle $[-5~;~3]$.
\end{enumerate}

\subsubsection*{Partie B}
On considère la fonction $f$ définie et dérivable sur l'intervalle $[-5~;~3]$ dont on donne ci-dessous la courbe représentative $C_f$.

\begin{center}
\begin{tikzpicture}
\begin{axis}[
    xmin=-5.3, xmax=3.5,
    ymin=-13, ymax=22,
    samples=200,
    domain=-5:3,
    axis lines=middle,
    axis line style={thick},
    xtick distance=0.5,
    ytick distance=5,
    tick label style={font=\scriptsize},
    x tick label style={
        /pgf/number format/.cd,
        fixed,
        fixed zerofill,
        precision=1,
        /tikz/.cd,
        /pgf/number format/zerofill=false
    },
    xticklabel={
        \pgfmathprintnumber{\tick}
    },
    x=1.5cm,
    grid=both,
    major grid style={gray!50, thick},
    minor grid style={gray!20, thin},
    minor x tick num=4,
    minor y tick num=4,
]
\addplot[red, very thick, restrict y to domain=-13:22] {(4*x^2 - 14*x + 8)*exp(0.5*x)};
\node[font=\scriptsize, below left] at (axis cs:0,0) {0};
\node[above left, color=red] at (axis cs:-4.5,16) {$C_f$};
\addplot[blue, mark=x, mark size=3pt] coordinates {(2,{(4*2^2 - 14*2 + 8)*exp(0.5*2)})};
\node[blue, below] at (axis cs:2,{(4*2^2 - 14*2 + 8)*exp(0.5*2)}) {A};
\pgfmathsetmacro{\yA}{(4*2^2 - 14*2 + 8)*exp(0.5*2)}
\addplot[blue, very thick] coordinates {(-5.3, \yA) (3.5, \yA)};
\node[blue, above] at (axis cs:-4.5,-11) {T};
\end{axis}
\end{tikzpicture}
\end{center}

La tangente $T$ à la courbe $C_f$ au point A d'abscisse 2 est horizontale.

\begin{enumerate}
\item Donner la valeur du nombre dérivé $f'(2)$.

\item Résoudre, avec la précision permise par le graphique, l'inéquation $f'(x) < 0$.

\item On sait que la fonction $f$ a pour expression sur l'intervalle $[-5~;~3]$ :
\[f(x) = (4x^2 - 14x + 8)\e^{0,5x}.\]

Démontrer que, pour tout $x$ appartenant à l'intervalle $[-5~;~3]$, on a :
\[f'(x) = P(x)\e^{0,5x}.\]

\item En utilisant les résultats de la \textbf{partie A}, dresser le tableau de variation de la fonction $f$ sur l'intervalle $[-5~;~3]$. \textit{(Il n'est pas demandé de calculer les images)}.
\end{enumerate}

\end{document}
